% Part: sets-functions-relations
% Chapter: relations
% Section: relations-as-sets

\documentclass[../../../include/open-logic-section]{subfiles}

\begin{document}

\olfileid{sfr}{rel}{set}
\olsection{ગણ તરીકે સંબંધો}

\begin{explain}
\olref[sfr][set][imp]{sec}માં આપણે કેટલાક મહત્વના ગણોનો ઉલ્લેખ કર્યો હતો:
$\Nat$, $\Int$, $\Rat$, $\Real$. આમાંના કેટલાક ગણોના !!{element}s
વચ્ચેના કેટલાક રસપ્રદ સંબંધો તમને ચોક્કસ યાદ હશે. દાખલા તરીકે,
આ દરેક ગણ પર એક સર્વસ્વીકૃત \emph{ક્રમ સંબંધ} છે.
\emph{તેની સાથે તાદાત્મ્ય ધરાવે છે} એવો સંબંધ પણ છે, જે દરેક વસ્તુ
પોતાની સાથે ધરાવે છે અને બીજી કોઈ વસ્તુ સાથે ધરાવતી નથી.
આગળ આપણે બીજા ઘણા રસપ્રદ સંબંધો જોઈશું, અને સંભવિત સંબંધો તો
એથીયે વધુ છે. જોકે, તેમની સમીક્ષા કરતાં પહેલાં આપણે એ નોંધથી
શરૂ કરીશું કે સંબંધોને ખાસ પ્રકારના ગણ તરીકે જોઈ શકાય છે.

આ માટે \olref[sfr][set][pai]{sec}માંથી બે બાબતો યાદ કરો. પ્રથમ,
\emph{ક્રમયુક્ત જોડ}નો ખ્યાલ યાદ કરો: $a$ અને $b$ આપેલાં હોય તો
આપણે $\tuple{a, b}$ રચી શકીએ છીએ. અહીં ઘટકોનો ક્રમ
\emph{મહત્વનો છે}. તેથી જો $a \neq b$, તો
$\tuple{a, b} \neq \tuple{b, a}$.
(આની સરખામણી અક્રમયુક્ત જોડ, એટલે કે $2$-ઘટકવાળા ગણ સાથે કરો,
જેમાં $\{a, b\}=\{b, a\}$.) બીજું, \emph{કાર્તેઝીય ગુણાકાર}નો
ખ્યાલ યાદ કરો: જો $A$ અને $B$ ગણો હોય, તો આપણે $A \times B$
રચી શકીએ છીએ; તે એવા બધા જોડ $\tuple{x, y}$નો ગણ છે જેમાં
$x \in A$ અને $y \in B$. ખાસ કરીને, $A^{2}= A \times A$
એ $A$માંથી મળતા બધા ક્રમયુક્ત જોડનો ગણ છે.

હવે આપણે એક ગણ પરનો ચોક્કસ સંબંધ લઈશું: પ્રાકૃતિક સંખ્યાઓના
ગણ $\Nat$ પરનો $<$-સંબંધ. સંખ્યાઓના એવા બધા જોડ
$\tuple{n, m}$નો ગણ લો જેમાં $n<m$, એટલે કે,
\[
  R=\Setabs{\tuple{n, m}}{n, m \in \Nat \text{ અને } n<m}.
\]
$n$ એ $m$ કરતાં નાનું હોવું અને જોડ $\tuple{n, m}$ એ
$R$નો સભ્ય હોવું, એ બે વચ્ચે ગાઢ જોડાણ છે, એટલે કે:
\[
      n<m\text{ તો અને તો જ }\tuple{n, m} \in R.
\]
ખરેખર, કોઈ માહિતી ગુમાવ્યા વિના, આપણે ગણ $R$ને
$\Nat$ પરનો $<$-સંબંધ \emph{જ છે} એમ ગણી શકીએ છીએ.

એ જ રીતે સંખ્યાઓ વચ્ચેના કોઈ પણ સંબંધ માટે આપણે $\Nat^{2}$નો
એક ઉપગણ રચી શકીએ છીએ. ઊલટું, સંખ્યાઓના જોડનો કોઈ પણ ગણ
$S \subseteq \Nat^{2}$ આપેલો હોય, તો તેને અનુરૂપ સંખ્યાઓ
વચ્ચે એક સંબંધ મળે છે, એટલે કે એવો સંબંધ જે $n$ એ $m$ સાથે
ધરાવે તો અને તો જ $\tuple{n, m} \in S$.
આ નીચેની વ્યાખ્યાને વાજબી ઠરાવે છે:
\end{explain}

\begin{defn}[દ્વિઘટકી સંબંધ]
ગણ $A$ પરનો \emph{દ્વિઘટકી સંબંધ} એ $A^{2}$નો ઉપગણ છે.
જો $R \subseteq A^{2}$ એ $A$ પરનો દ્વિઘટકી સંબંધ હોય અને
$x, y \in A$, તો આપણે ક્યારેક $\tuple{x, y} \in R$ માટે
$Rxy$ (અથવા $xRy$) લખીએ છીએ.
\end{defn}

\begin{ex}
  \ollabel{relations}
પ્રાકૃતિક સંખ્યાઓના જોડનો ગણ $\Nat^{2}$ નીચે પ્રમાણે
દ્વિપરિમાણીય શ્રેણિકમાં સૂચવી શકાય:
\[
  \begin{array}{ccccc}
  \mathbf{\tuple{ 0,0 }} & \tuple{ 0,1 } &
    \tuple{ 0,2 } & \tuple{ 0,3 } & \ldots\\
  \tuple{ 1,0 } & \mathbf{\tuple{ 1,1 }} &
    \tuple{ 1,2 } & \tuple{ 1,3 } & \ldots\\
  \tuple{ 2,0 } & \tuple{ 2,1 } &
    \mathbf{\tuple{ 2,2 }} & \tuple{ 2,3 } & \ldots\\
  \tuple{ 3,0 } & \tuple{ 3,1 } & \tuple{ 3,2 } &
    \mathbf{\tuple{ 3,3 }} & \ldots\\
  \vdots & \vdots & \vdots & \vdots & \mathbf{\ddots}
  \end{array}
\]
અહીં આપણે વિકર્ણ ઘાટા અક્ષરે દર્શાવ્યો છે, કારણ કે વિકર્ણ પર
આવેલા જોડનો બનેલો $\Nat^2$નો ઉપગણ, એટલે કે,
\[
  \{\tuple{0,0 }, \tuple{ 1,1 }, \tuple{ 2,2 }, \dots\},
  \]
એ $\Nat$ \emph{પરનો તાદાત્મ્ય સંબંધ} છે. (તાદાત્મ્ય સંબંધ
વારંવાર વપરાય છે, તેથી કોઈ પણ ગણ $A$ માટે
$\Id{A}=\Setabs{\tuple{ x,x }}{x \in A}$ એવી વ્યાખ્યા કરીએ.)
વિકર્ણની ઉપર આવેલા બધા જોડનો ઉપગણ, એટલે કે,
\[
  L = \{\tuple{ 0,1 },\tuple{ 0,2 },\ldots,\tuple{ 1,2 },
  \tuple{ 1,3 }, \dots, \tuple{ 2,3 }, \tuple{ 2,4 },\ldots\},
\]
એ \emph{કરતાં નાનું} એવો સંબંધ છે, એટલે કે $Lnm$ તો અને તો જ
$n<m$. વિકર્ણની નીચે આવેલા જોડનો ઉપગણ, એટલે કે,
\[
  G=\{ \tuple{ 1,0 },\tuple{ 2,0 },\tuple{
    2,1 }, \tuple{ 3,0 },\tuple{ 3,1 },\tuple{ 3,2 }, \dots\},
\]
એ \emph{કરતાં મોટું} એવો સંબંધ છે, એટલે કે $Gnm$ તો અને તો જ
$n>m$. $L$ અને $I$નો યોગગણ, જેને આપણે $K=L\cup I$ કહી શકીએ,
એ \emph{કરતાં નાનું અથવા સમાન} એવો સંબંધ છે:
$Knm$ તો અને તો જ $n \le m$. એ જ રીતે, $H=G \cup I$
એ \emph{કરતાં મોટું અથવા સમાન એવો સંબંધ} છે.
આ સંબંધો $L$, $G$, $K$ અને $H$ ખાસ પ્રકારના સંબંધો છે,
જેને \emph{ક્રમો} કહે છે. $L$ અને $G$નો એવો ગુણધર્મ છે કે
કોઈ સંખ્યા પોતાની સાથે $L$ કે $G$ ધરાવતી નથી
(એટલે કે, દરેક $n$ માટે, ન તો $Lnn$, ન તો $Gnn$).
આ ગુણધર્મવાળા સંબંધોને \emph{અસ્વવાચક} કહે છે અને,
જો તેઓ ક્રમો પણ હોય, તો તેમને \emph{કડક ક્રમો} કહે છે.
\end{ex}

\begin{explain}
ક્રમો અને તાદાત્મ્ય મહત્વના અને સ્વાભાવિક સંબંધો હોવા છતાં,
એ બાબત પર ભાર મૂકવો જોઈએ કે આપણી વ્યાખ્યા પ્રમાણે $A^{2}$નો
\emph{કોઈ પણ} ઉપગણ એ $A$ પરનો સંબંધ છે, ભલે તે ગમે તેટલો
અસ્વાભાવિક કે કૃત્રિમ લાગે. ખાસ કરીને, $\emptyset$ એ કોઈ પણ
ગણ પરનો સંબંધ છે (\emph{ખાલી સંબંધ}, જે ઘટકોનું કોઈ જોડ
ધરાવતું નથી), અને $A^{2}$ પોતે પણ $A$ પરનો સંબંધ છે
(જે દરેક જોડ ધરાવે છે), જેને \emph{સાર્વત્રિક સંબંધ} કહે છે.
પણ $E=\Setabs{\tuple{n, m}}{n>5 \text{ અથવા } m \times n \ge 34}$
જેવું કંઈક પણ સંબંધ ગણાય છે.
\end{explain}

\begin{prob}
  ગણ $\Pow{\{a, b, c\}}$ પરના સંબંધ $\subseteq$ના
  !!{element}sની યાદી આપો.
\end{prob}

\end{document}
